% Bagian: first-order-logic
% Bab: syntax-and-semantics
% Subbagian: satisfaction

\documentclass[../../../include/open-logic-section]{subfiles}

\begin{document}

\olfileid[id]{fol}{syn}{sat}

\olsection{Keterpenuhan \article{formula} \printtoken{S}{formula}
  dalam \article{structure} \printtoken{S}{structure}}

\begin{explain}
Gagasan-gagasan dasar yang menghubungkan ungkapan seperti suku dan
!!{formula}s, di satu pihak, dengan !!{structure}s, di pihak lain, ialah
\emph{!!{value}} suatu suku dan \emph{keterpenuhan} !!a{formula}. Secara
informal, !!{value} suatu suku merupakan !!{element} dari
!!a{structure}---jika suku itu hanya berupa konstanta, !!{value}-nya ialah
objek yang ditetapkan kepada konstanta tersebut oleh !!{structure}, dan jika
suku itu dibangun dengan menggunakan !!{function}s, !!{value}-nya dihitung
dari !!{value}s konstanta dan fungsi-fungsi yang ditetapkan kepada
simbol fungsi dalam suku tersebut. !!^a{formula} \emph{terpenuhi} dalam
!!a{structure} jika interpretasi yang diberikan kepada predikat membuat
!!{formula} tersebut benar dalam domain !!{structure}. Gagasan keterpenuhan
ini ditentukan secara induktif: spesifikasi !!{structure} secara langsung
menyatakan kapan !!{formula}s atomik terpenuhi, dan kita mendefinisikan kapan
suatu !!{formula} kompleks terpenuhi berdasarkan penghubung utama atau
kuantornya dan berdasarkan apakah !!{subformula}s langsungnya terpenuhi atau
tidak.

Kasus kuantor di sini agak rumit, sebab !!{subformula} langsung dari suatu
!!{formula} berkuantor dapat mempunyai !!{variable} bebas, sedangkan
!!{structure}s tidak menentukan !!{value}s dari !!{variable}s. Untuk
mengatasi kesulitan ini, kita juga memperkenalkan \emph{penugasan variabel}
dan mendefinisikan keterpenuhan bukan hanya sehubungan dengan
!!a{structure}, melainkan sehubungan dengan !!a{structure} beserta
penugasan !!a{variable}.
\end{explain}

\begin{defn}[Penugasan Variabel]
Suatu \emph{penugasan variabel}~$s$ bagi !!a{structure}~$\Struct{M}$ ialah
fungsi yang memetakan setiap !!{variable} ke !!a{element} dari~$\Domain M$,
yakni, $s\colon \Var \to \Domain M$.
\end{defn}

\begin{explain}
!!^a{structure} menetapkan !!a{value} kepada setiap !!{constant}, dan
penugasan variabel menetapkan nilai kepada setiap variabel. Namun, kita juga
ingin menggunakan suku yang dibangun dari unsur-unsur tersebut untuk menamai
!!{element}s dari !!{domain}. Untuk itu, kita mendefinisikan !!{value} suku
secara induktif. Bagi !!{constant}s dan variabel, nilainya tepat seperti yang
ditentukan oleh !!{structure} atau penugasan variabel; bagi suku yang lebih
kompleks, nilainya dihitung secara rekursif dengan menggunakan fungsi-fungsi
yang ditetapkan oleh !!{structure} kepada !!{function}s.
\end{explain}

\begin{defn}[!!^{value} Suku]
Jika $t$ merupakan suku dari bahasa~$\Lang L$, $\Struct M$ merupakan
!!{structure} bagi~$\Lang L$, dan $s$ merupakan penugasan
!!a{variable} bagi~$\Struct M$, maka \emph{!!{value}}~$\Value{t}{M}[s]$
didefinisikan sebagai berikut:
\begin{enumerate}
\item \indcase{t}{c}{$\Value{\indfrm}{M}[s] = \Assign{\indcomplex}{M}$.}
\item \indcase{t}{x}{$\Value{\indfrm}{M}[s] = s(\indcomplex)$.}
\item \indcase{t}{\Atom{f}{t_1, \ldots, t_n}}{
\[
\Value{\indfrm}{M}[s] = \Assign{f}{M}(\Value{t_1}{M}[s], \ldots,
\Value{t_n}{M}[s]).
\]}
\end{enumerate}
\end{defn}

\begin{defn}[Varian-$x$]
Jika $s$ merupakan penugasan !!a{variable} bagi !!a{structure}~$\Struct M$,
maka setiap penugasan !!{variable}~$s'$ bagi~$\Struct M$ yang hanya mungkin
berbeda dari~$s$ dalam nilai yang ditetapkannya kepada~$x$ disebut
\emph{varian-$x$} dari~$s$. Jika $s'$ merupakan varian-$x$ dari~$s$, kita
menulis $\varAssign{s'}{s}{x}$.
\end{defn}

\begin{explain}
Perhatikan bahwa varian-$x$ dari suatu penugasan~$s$ tidak
\emph{harus} menetapkan sesuatu yang berbeda kepada~$x$. Bahkan, setiap
penugasan terhitung sebagai varian-$x$ dari dirinya sendiri.
\end{explain}

\begin{defn}
  Jika $s$ merupakan penugasan !!a{variable} bagi !!a{structure}~$\Struct M$
  dan $m \in \Domain{M}$, maka penugasan~$\Subst{s}{m}{x}$ ialah penugasan
  variabel yang didefinisikan oleh
  \[\Subst{s}{m}{x}(y) = \begin{cases}
    m & \text{jika } y \ident x\\
    s(y) & \text{jika tidak}.
  \end{cases}\]
\end{defn}

Dengan kata lain, $\Subst{s}{m}{x}$ merupakan varian-$x$ tertentu dari~$s$
yang menetapkan !!{element} domain~$m$ kepada~$x$, dan kepada
!!{variable}s selain~$x$ menetapkan hal yang sama seperti yang
ditetapkan~$s$.

\begin{defn}[Keterpenuhan]
\ollabel{defn:satisfaction}
Keterpenuhan !!a{formula}~$!A$ dalam !!a{structure}~$\Struct M$
relatif terhadap penugasan !!a{variable}~$s$, yang secara simbolis ditulis
$\Sat{M}{!A}[s]$, didefinisikan secara rekursif sebagai berikut. (Kita
menulis $\Sat/{M}{!A}[s]$ untuk berarti ``tidak $\Sat{M}{!A}[s]$.'')
\begin{enumerate}
\tagitem{prvFalse}{%
  \indcase{!A}{\lfalse}{$\Sat/{M}{\indfrm}[s]$.}}{}

\tagitem{prvTrue}{%
  \indcase{!A}{\ltrue}{$\Sat{M}{\indfrm}[s]$.}}{}

\item \indcase{!A}{\Atom{R}{t_1, \dots, t_n}}{$\Sat{M}{\indfrm}[s]$
  jika dan hanya jika $\langle \Value{t_1}{M}[s], \dots,
  \Value{t_n}{M}[s] \rangle \in \Assign{R}{M}$.}

\item \indcase{!A}{\eq[t_1][t_2]}{$\Sat{M}{\indfrm}[s]$ jika dan hanya jika
  $\Value{t_1}{M}[s] = \Value{t_2}{M}[s]$.}

\tagitem{prvNot}{%
  \indcase{!A}{\lnot !B}{$\Sat{M}{\indfrm}[s]$ jika dan hanya jika
    $\Sat/{M}{!B}[s]$.}}{}

\tagitem{prvAnd}{%
  \indcase{!A}{(!B \land !C)}{$\Sat{M}{\indfrm}[s]$ jika dan hanya jika
    $\Sat{M}{!B}[s]$ dan $\Sat{M}{!C}[s]$.}}{}

\tagitem{prvOr}{%
  \indcase{!A}{(!B \lor !C)}{$\Sat{M}{\indfrm}[s]$ jika dan hanya jika
    $\Sat{M}{!B}[s]$ atau $\Sat{M}{!C}[s]$ (atau keduanya).}}{}

\tagitem{prvIf}{%
  \indcase{!A}{(!B \lif !C)}{$\Sat{M}{\indfrm}[s]$ jika dan hanya jika
    $\Sat/{M}{!B}[s]$ atau $\Sat{M}{!C}[s]$ (atau keduanya).}}{}

\tagitem{prvIff}{%
  \indcase{!A}{(!B \liff !C)}{$\Sat{M}{\indfrm}[s]$ jika dan hanya jika
    baik $\Sat{M}{!B}[s]$ maupun $\Sat{M}{!C}[s]$ keduanya berlaku, atau
    baik $\Sat{M}{!B}[s]$ maupun $\Sat{M}{!C}[s]$ keduanya tidak berlaku.}}{}

\tagitem{prvAll}{%
  \indcase{!A}{\lforall[x][!B]}{$\Sat{M}{\indfrm}[s]$ jika dan hanya jika
    bagi setiap !!{element}~$m \in \Domain M$,
    $\Sat{M}{!B}[\Subst{s}{m}{x}]$.}}{}

\tagitem{prvEx}{%
  \indcase{!A}{\lexists[x][!B]}{$\Sat{M}{\indfrm}[s]$ jika dan hanya jika
  bagi setidaknya satu !!{element}~$m \in \Domain M$,
  $\Sat{M}{!B}[\Subst{s}{m}{x}]$.}}{}
\end{enumerate}
\end{defn}

\begin{explain}
Penugasan variabel penting dalam
\iftag{notprvEx,notprvAll}{klausa terakhir}{dua klausa terakhir}.\iftag{prvAll}{
Kita tidak dapat mendefinisikan keterpenuhan $\lforall[x][!B(x)]$ dengan
``bagi semua $m \in \Domain{M}$, $\Sat{M}{!B(m)}$.''}{}\iftag{prvEx}{ Kita
tidak dapat mendefinisikan keterpenuhan $\lexists[x][!B(x)]$ dengan ``bagi
setidaknya satu $m \in \Domain{M}$, $\Sat{M}{!B(m)}$.''}{} Alasannya ialah
bahwa jika $m \in \Domain M$, anggota itu bukan simbol bahasa, sehingga
$!B(m)$~bukan !!a{formula} (yakni, $\Subst{!B}{m}{x}$ tidak terdefinisi).
Kita juga tidak dapat mengandaikan bahwa tersedia !!{constant}s atau suku
yang menamai setiap !!{element} dari~$\Struct{M}$, sebab tidak ada apa pun
dalam definisi !!{structure}s yang mengharuskannya. Dalam bahasa standar,
himpunan !!{constant}s bersifat !!{denumerable}, sehingga jika
$\Domain{M}$ tidak !!{enumerable}, bahkan tidak tersedia cukup
!!{constant}s untuk menamai setiap objek.

Kita menyelesaikan masalah ini dengan memperkenalkan penugasan
!!{variable}, yang memungkinkan kita menghubungkan variabel secara langsung
dengan !!{element}s domain. Kemudian, alih-alih mengatakan bahwa, misalnya,
\iftag{prvEx}{$\lexists[x][!B(x)]$}{$\lforall[x][!B(x)]$} terpenuhi
dalam~$\Struct M$ jika dan hanya jika \iftag{prvEx}{bagi setidaknya satu
$m \in \Domain{M}$, $\Sat{M}{!B(m)}$}{bagi semua $m \in \Domain{M}$,
$\Sat{M}{!B(m)}$}, kita mengatakan bahwa formula itu terpenuhi
dalam~$\Struct M$ \emph{relatif terhadap}~$s$ jika dan hanya jika $!B(x)$
terpenuhi relatif terhadap~$\Subst{s}{m}{x}$
\iftag{prvEx}{bagi setidaknya satu}{bagi setiap} $m \in \Domain M$.
\end{explain}

\begin{ex}
Misalkan $\Lang{L} = \{a, b, f, R\}$, dengan $a$ dan $b$ merupakan
!!{constant}s, $f$~merupakan !!{function} dua-tempat, dan $R$~merupakan
!!{predicate} dua-tempat. Perhatikan !!{structure}~$\Struct{M}$ yang
didefinisikan oleh:
\begin{enumerate}
\item $\Domain M = \{1, 2, 3, 4\}$
\item $\Assign{a}{M} = 1$
\item $\Assign{b}{M} = 2$
\item $\Assign{f}{M}(x, y) = x+y$ jika $x+y \le 3$ dan $= 3$ jika tidak.
\item $\Assign{R}{M} = \{\tuple{1, 1}, \tuple{1, 2}, \tuple{2, 3}, \tuple{2, 4}\}$
\end{enumerate}
Fungsi $s(x) = 1$ yang menetapkan $1 \in \Domain{M}$ kepada setiap
!!{variable} merupakan penugasan variabel bagi~$\Struct{M}$.

Maka
\begin{align*}
\Value{f(a,b)}{M}[s] & = \Assign{f}{M}(\Value{a}{M}[s], \Value{b}{M}[s]).
\intertext{Karena $a$ dan $b$ merupakan !!{constant}s, $\Value{a}{M}[s]
  = \Assign{a}{M} = 1$ dan $\Value{b}{M}[s] = \Assign{b}{M} = 2$. Jadi}
\Value{f(a,b)}{M}[s] & = \Assign{f}{M}(1, 2) = 1+2 = 3.
\intertext{Untuk menghitung nilai $f(f(a,b),a)$, kita harus memperhatikan}
\Value{f(f(a,b),a)}{M}[s] & = \Assign{f}{M}(\Value{f(a, b)}{M}[s],
\Value{a}{M}[s]) = \Assign{f}{M}(3, 1) = 3,
\intertext{karena $3+1 > 3$. Karena $s(x) = 1$ dan $\Value{x}{M}[s] =
  s(x)$, kita juga mempunyai}
\Value{f(f(a,b),x)}{M}[s] & = \Assign{f}{M}(\Value{f(a, b)}{M}[s],
\Value{x}{M}[s]) = \Assign{f}{M}(3, 1) = 3,
\end{align*}

Suatu !!{formula} atomik~$R(t_1, t_2)$ terpenuhi jika tupel nilai
argumen-argumennya, yakni, $\tuple{\Value{t_1}{M}[s],
  \Value{t_2}{M}[s]}$, merupakan !!a{element} dari~$\Assign{R}{M}$. Jadi,
misalnya, kita mempunyai $\Sat{M}{R(b,f(a,b))}[s]$ karena
$\tuple{\Value{b}{M}, \Value{f(a,b)}{M}} = \tuple{2, 3} \in
  \Assign{R}{M}$, tetapi $\Sat/{M}{R(x, f(a,b))}[s]$ karena
  $\tuple{1, 3} \notin \Assign{R}{M}$.

Untuk menentukan apakah suatu formula nonatomik~$!A$ terpenuhi, Anda
menerapkan klausa dalam definisi induktif yang berlaku bagi penghubung
utama. Misalnya, penghubung utama dalam $R(a, a) \lif (R(b, x) \lor R(x,
b))$ ialah~$\lif$, dan
\begin{align*}
  & \Sat{M}{R(a, a) \lif (R(b, x) \lor R(x, b))}[s]
    \text{ jika dan hanya jika }\\
  & \qquad
  \Sat/{M}{R(a,a)}[s] \text{ atau } \Sat{M}{R(b, x) \lor R(x, b)}[s]
  \intertext{Karena $\Sat{M}{R(a,a)}[s]$ (sebab $\tuple{1,1} \in
    \Assign{R}{M}$), kita belum dapat menentukan jawabannya dan pertama-tama
    harus mencari tahu apakah $\Sat{M}{R(b, x) \lor R(x, b)}[s]$:}
  & \Sat{M}{R(b, x) \lor R(x, b)}[s] \text{ jika dan hanya jika }\\
  & \qquad \Sat{M}{R(b,x)}[s] \text{ atau } \Sat{M}{R(x, b)}[s]
  \intertext{Dan memang demikian, sebab $\Sat{M}{R(x, b)}[s]$
    (karena $\tuple{1,2} \in \Assign{R}{M}$).}
\end{align*}

Ingat bahwa varian-$x$ dari~$s$ merupakan penugasan variabel yang
berbeda dari $s$ paling banyak dalam apa yang ditetapkannya kepada~$x$.
Bagi setiap !!{element} dari~$\Domain{M}$, terdapat varian-$x$ dari~$s$:
\begin{align*}
  s_1 & = \Subst{s}{1}{x}, &
  s_2 & = \Subst{s}{2}{x},\\
  s_3 & = \Subst{s}{3}{x}, &
  s_4 & = \Subst{s}{4}{x}.
\end{align*}
Jadi, misalnya, $s_2(x) = 2$ dan $s_2(y) = s(y) = 1$ bagi semua
variabel~$y$ selain~$x$. Inilah semua varian-$x$ dari~$s$ bagi
struktur~$\Struct{M}$, sebab $\Domain{M} = \{1, 2, 3, 4\}$. Secara khusus,
perhatikan bahwa $s_1 = s$ ($s$~selalu merupakan varian-$x$ dari dirinya
sendiri).

\iftag{prvEx}{Untuk menentukan apakah suatu !!{formula} berkuantor
  eksistensial~$\lexists[x][!A(x)]$ terpenuhi, kita harus menentukan apakah
  $\Sat{M}{!A(x)}[\Subst{s}{m}{x}]$ bagi setidaknya satu $m \in \Domain M$.
  Jadi,
  \[
  \Sat{M}{\lexists[x][(R(b,x) \lor R(x,b))]}[s],
  \]
  sebab $\Sat{M}{R(b,x) \lor R(x, b)}[\Subst{s}{1}{x}]$
  ($\Subst{s}{3}{x}$ juga memenuhi keperluan). Namun,
  \[
  \Sat/{M}{\lexists[x][(R(b,x) \land R(x,b))]}[s]
  \]
  sebab, apa pun $m \in \Domain{M}$ yang kita pilih,
  $\Sat/{M}{R(b,x) \land R(x,b)}[\Subst{s}{m}{x}]$.}{}

\iftag{prvAll}{Untuk menentukan apakah suatu !!{formula} berkuantor
  universal~$\lforall[x][!A(x)]$ terpenuhi, kita harus menentukan apakah
  $\Sat{M}{!A(x)}[\Subst{s}{m}{x}]$ bagi semua $m \in \Domain M$. Jadi,
  \[
  \Sat{M}{\lforall[x][(R(x,a) \lif R(a,x))]}[s],
  \]
  sebab $\Sat{M}{R(x,a) \lif R(a,x)}[\Subst{s}{m}{x}]$ bagi semua
  $m \in \Domain M$. Untuk $m = 1$, kita mempunyai
  $\Sat{M}{R(a,x)}[\Subst{s}{1}{x}]$, sehingga konsekuennya benar; untuk
  $m = 2$, $3$, dan~$4$, kita mempunyai
  $\Sat/{M}{R(x,a)}[\Subst{s}{m}{x}]$, sehingga antesedennya salah. Namun,
  \[
  \Sat/{M}{\lforall[x][(R(a,x) \lif R(x,a))]}[s]
  \]
  sebab $\Sat/{M}{R(a,x) \lif R(x,a)}[\Subst{s}{2}{x}]$ (karena
  $\Sat{M}{R(a, x)}[\Subst{s}{2}{x}]$ dan $\Sat/{M}{R(x,
  a)}[\Subst{s}{2}{x}]$).}{}

\iftag{defEx}{Untuk menentukan apakah suatu !!{formula} berkuantor
  eksistensial~$\lexists[x][!A(x)]$ terpenuhi, kita harus menentukan apakah
  $\Sat{M}{\lnot \lforall[x][\lnot !A(x)]}[s]$. Misalnya, kita mempunyai
  \[
  \Sat{M}{\lexists[x][(R(b,x) \lor R(x,b))]}[s].
  \]
  Pertama, $\Sat{M}{R(b,x) \lor R(x, b)}[\Subst{s}{1}{x}]$
  ($\Subst{s}{3}{x}$ juga memenuhi keperluan). Jadi,
  $\Sat/{M}{\lnot(R(b,x) \lor R(x,b))}[\Subst{s}{1}{x}]$, dengan demikian
  $\Sat/{M}{\lforall[x][\lnot(R(b,x) \lor R(x,b))]}[s]$, dan
  oleh karena itu $\Sat{M}{\lnot\lforall[x][\lnot(R(b,x) \lor
  R(x,b))]}[s]$. Di lain pihak,
  \[
  \Sat/{M}{\lexists[x][(R(b,x) \land R(x,b))]}[s].
  \]
  Alasannya ialah $\Sat{M}{\lforall[x][\lnot(R(b,x) \land
  R(x,b))]}[s]$, sebab untuk tidak satu pun $m \in \Domain M$ berlaku
  $\Sat{M}{R(b,x) \land R(x,b)}[\Subst{s}{m}{x}]$. Seperti yang mungkin
  dapat Anda duga dari contoh-contoh ini,
  $\Sat{M}{\lexists[x][!A(x)]}[s]$ jika dan hanya jika
  $\Sat{M}{!A(x)}[\Subst{s}{m}{x}]$ bagi setidaknya satu
  $m \in \Domain M$.}{}

\iftag{defAll}{Untuk menentukan apakah suatu !!{formula} berkuantor
  universal~$\lforall[x][!A(x)]$ terpenuhi, kita harus menentukan apakah
  $\Sat{M}{\lnot\lexists[x][\lnot !A(x)]}[s]$. Misalnya,
  \[
  \Sat{M}{\lforall[x][(R(x,a) \lif R(a,x))]}[s],
  \]
  Pertama, $\Sat{M}{R(x,a) \lif R(a,x)}[\Subst{s}{m}{x}]$ bagi semua
  $m \in \Domain M$ ($\Sat{M}{R(a,x)}[\Subst{s}{1}{x}]$ dan
  $\Sat/{M}{R(x,a)}[\Subst{s}{m}{x}]$ untuk $m = 2$, $3$, atau~$4$).
  Jadi, tidak ada $m \in \Domain M$ sedemikian sehingga
  $\Sat{M}{\lnot(R(x,a) \lif R(a,x))}[\Subst{s}{m}{x}]$, dan karenanya
  $\Sat/{M}{\lexists[x][\lnot(R(x,a) \lif R(a,x))]}[s]$. Oleh karena itu,
  $\Sat{M}{\lnot\lexists[x][\lnot(R(x,a) \lif R(a,x))]}[s]$. Di lain pihak,
  \[
  \Sat/{M}{\lforall[x][(R(a,x) \lif R(x,a))]}[s],
  \]
  sebab $\Sat/{M}{R(a,x) \lif R(x,a)}[\Subst{s}{2}{x}]$, dan dengan demikian
  $\Sat{M}{\lexists[x][\lnot(R(a,x) \lif R(x,a))]}[s]$. Seperti yang mungkin
  dapat Anda duga dari contoh-contoh ini,
  $\Sat{M}{\lforall[x][!A(x)]}[s]$ jika dan hanya jika
  $\Sat{M}{!A(x)}[\Subst{s}{m}{x}]$ bagi setiap $m \in \Domain M$.}{}

Untuk kasus yang lebih rumit, perhatikan
\[
\lforall[x][(R(a,x) \lif \lexists[y][R(x,y)])].
\]
Karena $\Sat/{M}{R(a,x)}[\Subst{s}{3}{x}]$ dan
$\Sat/{M}{R(a,x)}[\Subst{s}{4}{x}]$, kasus menarik yang di dalamnya kita
harus memperhatikan konsekuen dari kondisional itu hanya $m = 1$ dan
$m = 2$. Apakah $\Sat{M}{\lexists[y][R(x,y)]}[\Subst{s}{1}{x}]$ berlaku?
Hal itu berlaku jika terdapat setidaknya satu $n \in \Domain M$ sedemikian
sehingga $\Sat{M}{R(x,y)}[\Subst{\Subst{s}{1}{x}}{n}{y}]$. Bahkan, jika
kita mengambil $n = 1$, kita mempunyai
$\Subst{\Subst{s}{1}{x}}{n}{y} = \Subst{s}{1}{y} = s$. Karena
$s(x) = 1$, $s(y) = 1$, dan $\tuple{1,1} \in \Assign{R}{M}$, jawabannya ya.

Untuk menentukan apakah $\Sat{M}{\lexists[y][R(x,y)]}[\Subst{s}{2}{x}]$,
kita harus memperhatikan penugasan !!{variable}
$\Subst{\Subst{s}{2}{x}}{n}{y}$. Di sini, untuk $n = 1$, penugasan ini
ialah~$s_2 = \Subst{s}{2}{x}$, yang tidak memenuhi $R(x,y)$ ($s_2(x) = 2$,
$s_2(y) = 1$, dan $\tuple{2,1}\notin \Assign{R}{M}$). Namun, perhatikan
$\Subst{\Subst{s}{2}{x}}{3}{y} = \Subst{s_2}{3}{y}$.
$\Sat{M}{R(x,y)}[\Subst{s_2}{3}{y}]$ sebab $\tuple{2,3} \in
\Assign{R}{M}$, sehingga $\Sat{M}{\lexists[y][R(x,y)]}[s_2]$.

Jadi, bagi semua $m \in \Domain M$, berlaku
$\Sat/{M}{R(a,x)}[\Subst{s}{m}{x}]$ (jika $m = 3$, $4$) atau
$\Sat{M}{\lexists[y][R(x,y)]}[\Subst{s}{m}{x}]$ (jika $m = 1$, $2$),
sehingga
\[
\Sat{M}{\lforall[x][(R(a,x) \lif \lexists[y][R(x,y)])]}[s].
\]
Di lain pihak,
\[
\Sat/{M}{\lexists[x][(R(a,x) \land \lforall[y][R(x,y)])]}[s].
\]
Kita mempunyai $\Sat{M}{R(a,x)}[\Subst{s}{m}{x}]$ hanya untuk $m = 1$ dan
$m = 2$. Namun, untuk kedua nilai~$m$ ini, pada gilirannya terdapat suatu
$n \in \Domain M$---yakni $n = 4$ ketika $m = 1$, dan $n = 1$ ketika
$m = 2$---sedemikian sehingga
$\Sat/{M}{R(x,y)}[\Subst{\Subst{s}{m}{x}}{n}{y}]$, dan dengan demikian
$\Sat/{M}{\lforall[y][R(x,y)]}[\Subst{s}{m}{x}]$ untuk $m = 1$ dan
$m = 2$. Singkatnya, tidak ada $m \in \Domain M$ sedemikian sehingga
$\Sat{M}{R(a,x) \land \lforall[y][R(x,y)]}[\Subst{s}{m}{x}]$.
\end{ex}

\iftag{defEx}{%
\begin{prop}\ollabel{prop:sat-ex}
$\Sat{M}{\lexists[x][!B(x)]}[s]$ jika dan hanya jika terdapat varian-$x$
$s'$ dari $s$ sedemikian sehingga $\Sat{M}{!B(x)}[s']$.
\end{prop}

\begin{proof}
  Latihan.
\end{proof}
}{}

\tagprob{defEx}
\begin{prob}
  Buktikan \olref[fol][syn][sat]{prop:sat-ex}
\end{prob}
\tagendprob

\iftag{defAll}{%
\begin{prop}\ollabel{prop:sat-all}
$\Sat{M}{\lforall[x][!B(x)]}[s]$ jika dan hanya jika bagi setiap
varian-$x$~$s'$ dari $s$, $\Sat{M}{!B(x)}[s']$
\end{prop}

\begin{proof}
  Latihan.
\end{proof}
}{}

\tagprob{defAll}
\begin{prob}
  Buktikan \olref[fol][syn][sat]{prop:sat-all}
\end{prob}
\tagendprob

\begin{prob}
Misalkan $\Lang L = \{c, f, A\}$ dengan satu !!{constant}, satu
!!{function} satu-tempat, dan satu !!{predicate} dua-tempat, serta misalkan
!!{structure}~$\Struct{M}$ diberikan oleh
\begin{enumerate}
\item $\Domain M = \{1, 2, 3\}$
\item $\Assign{c}{M} = 3$
\item $\Assign{f}{M}(1) = 2, \Assign{f}{M}(2) = 3, \Assign{f}{M}(3) = 2$
\item $\Assign{A}{M} = \{\tuple{1, 2}, \tuple{2, 3}, \tuple{3, 3}\}$
\end{enumerate}
(a) Misalkan $s(v) = 1$ bagi semua !!{variable}s~$v$. Tentukan apakah
\[
\Sat{M}{\lexists[x][(A(f(z), c) \lif \lforall[y][(A(y, x) \lor A(f(y),
      x))])]}[s]
\]
berlaku. Jelaskan mengapa berlaku atau mengapa tidak.

(b) Berikan struktur dan penugasan !!{variable} lain yang di dalamnya
!!{formula} tersebut tidak terpenuhi.
\end{prob}

\end{document}
